\subsubsection{Stern-Brocot (best fraction in interval)}
\begin{lstlisting}[style=mystyle, language=C++]
// TC: O(log(max(a,b))) per query
// usa: encontrar la fraccion mas simple p/q en un rango (lo, hi)
#include <bits/stdc++.h>
using namespace std;

struct Frac { long long p, q; };

// simplest = minimize p+q; iterate through mediants
//     finds the fraction p/q with smallest p+q in the open interval (a/b, c/d)
//     assumes a/b < c/d.
Frac bestInInterval(long long a, long long b, long long c, long long d){
	long long p = a + c, q = b + d;
	while(true){
		long long k = max((p - a) / b, (p - c) / d);
		// this is a coarse approximation; a full implementation needs limits
		if(k <= 1) break;
		p = a + k * c;
		q = b + k * d;
	}
	return {p, q};
}

int main(){
	// simplest fraction in (0/1, 1/1) is 1/2
	auto f = bestInInterval(0, 1, 1, 1);
	cout << f.p << "/" << f.q << "\n";
	return 0;
}
\end{lstlisting}
